\documentclass[11pt,a4paper,UTF8]{ctexart}
\usepackage{geometry}
\geometry{left=18mm,right=18mm,top=24mm,bottom=24mm,headheight=22pt,footskip=12mm}
\usepackage{amsmath,amssymb,mathtools,bm,esint,extarrows}
\usepackage{xcolor}
\usepackage{tcolorbox}
\tcbuselibrary{skins,breakable}
\usepackage{fancyhdr}
\usepackage{titlesec}
\usepackage{hyperref}
\usepackage{tikz}
\usepackage{booktabs}
\usepackage{tabularx}

% --- Color Palette (Purple & DeepPink Academic Theme) ---
\definecolor{DeepPurple}{HTML}{4C1D95}
\definecolor{TitlePurple}{HTML}{581C87}
\definecolor{SubPurple}{HTML}{6B21A8}
\definecolor{DeepPink}{HTML}{FF1493}
\definecolor{LightPinkBg}{HTML}{FFF5F8}
\definecolor{LightPurpleBg}{HTML}{F8F5FF}
\definecolor{GoldAccent}{HTML}{D97706}
\definecolor{DarkText}{HTML}{1F2937}

\hypersetup{
    colorlinks=true,
    linkcolor=TitlePurple,
    citecolor=DeepPink,
    urlcolor=SubPurple,
    pdfborder={0 0 0}
}

% --- Typography & Spacing ---
\linespread{1.3}
\setlength{\parskip}{0.35em plus 0.1em minus 0.1em}
\setlength{\parindent}{0pt}

% --- Section Titles Styling ---
\titleformat{\section}
  {\Large\bfseries\color{TitlePurple}}{\thesection}{1em}{}
  [\vspace{1mm}{\color{TitlePurple!30}\hrule height 1pt}]
\titleformat{\subsection}
  {\large\bfseries\color{SubPurple}}{\thesubsection}{0.8em}{}
\titleformat{\subsubsection}
  {\normalsize\bfseries\color{DeepPurple}}{\thesubsubsection}{0.6em}{}

% --- tcolorbox Environments ---
% 1. Theorem / Formula / Method Box (Deep Purple)
\newtcolorbox{thmbox}[1][]{
    enhanced,
    breakable,
    colback=LightPurpleBg,
    colframe=TitlePurple,
    arc=3pt,
    boxrule=0pt,
    leftrule=3.5pt,
    left=10pt,
    right=10pt,
    top=8pt,
    bottom=8pt,
    title=#1,
    coltitle=white,
    colbacktitle=TitlePurple,
    fonttitle=\bfseries\small,
    attach boxed title to top left={yshift=-2mm, xshift=4mm},
    boxed title style={boxrule=0pt, arc=2pt}
}

% 2. Solution / Formula / Derivation Box (DeepPink)
\newtcolorbox{solbox}[1][]{
    enhanced,
    breakable,
    colback=LightPinkBg,
    colframe=DeepPink,
    arc=3pt,
    boxrule=0pt,
    leftrule=3.5pt,
    left=10pt,
    right=10pt,
    top=8pt,
    bottom=8pt,
    title=#1,
    coltitle=white,
    colbacktitle=DeepPink,
    fonttitle=\bfseries\small,
    attach boxed title to top left={yshift=-2mm, xshift=4mm},
    boxed title style={boxrule=0pt, arc=2pt}
}

% 3. Learning Goals / Summary Box
\newtcolorbox{targetbox}[1][]{
    enhanced,
    breakable,
    colback=LightPurpleBg,
    colframe=DeepPurple,
    arc=3pt,
    boxrule=1pt,
    left=10pt,
    right=10pt,
    top=8pt,
    bottom=8pt,
    title=#1,
    coltitle=white,
    colbacktitle=DeepPurple,
    fonttitle=\bfseries\small,
    attach boxed title to top left={yshift=-2mm, xshift=4mm},
    boxed title style={boxrule=0pt, arc=2pt}
}

\begin{document}

% ------------------- 讲义封面 (Cover Page) -------------------
\begin{titlepage}
\thispagestyle{empty}
\begin{tikzpicture}[remember picture,overlay]
    \fill[DeepPurple] (current page.north west) rectangle ([yshift=-7cm]current page.north east);
    \draw[DeepPink, line width=3.5pt] ([yshift=-7cm]current page.north west) -- ([yshift=-7cm]current page.north east);
    \draw[GoldAccent, line width=1pt] ([yshift=-7.12cm]current page.north west) -- ([yshift=-7.12cm]current page.north east);
    
    \node[anchor=north east, opacity=0.12, text=white] at ([xshift=-1.5cm, yshift=-0.5cm]current page.north east) {
        \fontsize{70}{70}\selectfont\bfseries 2027
    };
\end{tikzpicture}

\vspace*{0.8cm}
\begin{center}
    {\color{white}\fontsize{26}{32}\selectfont\bfseries 2027 考研数学(二) 核心讲义}\\[0.6cm]
    {\color{white}\fontsize{13}{18}\selectfont\bfseries 全国硕士研究生招生考试 · 统考数学(二) 核心精讲全书}\\[1.5cm]
\end{center}

\vspace{3.5cm}
\begin{center}
    {\color{GoldAccent}\large\bfseries $\blacktriangleright$ 基础强化 · 核心题解 · 考点深水区 $\blacktriangleleft$}\\[0.8cm]
    {\color{TitlePurple}\fontsize{24}{28}\selectfont\bfseries [强化]一元积分学的应用(几何)}\\[0.6cm]
    {\color{DeepPink}\large\bfseries —— 考研高阶冲刺与深水区专攻 ——}\\[1.5cm]
\end{center}

\vfill

\begin{center}
\begin{tcolorbox}[
    colback=white,
    colframe=DeepPurple,
    width=0.88\textwidth,
    arc=4pt,
    boxrule=1.2pt,
    halign=center,
    drop shadow=gray!30
]
    \vspace{3mm}
    {\bfseries\large 编著团队：EscoffierZhou}\\[2.5mm]
    {\bfseries\color{TitlePurple} 目标院校：中国科学院大学 (UCAS) 计算机专硕 (22408)}\\[2.5mm]
    {\small\color{gray} 备考铁律：手算真题数字 · 错题白纸重做 · 408 客观题为王}\\[2.5mm]
    {\small 编制版本：2027 届首发版 · \today}
    \vspace{2mm}
\end{tcolorbox}
\end{center}
\vspace{0.8cm}
\end{titlepage}

% ------------------- 目录与页面主体配置 -------------------
\pagestyle{fancy}
\fancyhf{}
\fancyhead[L]{\small\color{gray}2027 考研数学(二) · 核心讲义系列}
\fancyhead[R]{\small\color{TitlePurple}\nouppercase{\leftmark}}
\fancyfoot[C]{\small\color{gray}— 第 \thepage\ 页 —}
\fancyfoot[R]{\small\color{gray}UCAS 22408}
\renewcommand{\headrulewidth}{0.5pt}

\pagenumbering{Roman}
\tableofcontents
\newpage
\pagenumbering{arabic}


\begin{targetbox}[本章复习目标与核心知识图谱]
\begin{itemize}
  \item 目的[1]:攻克平面任意倾斜直线轴 $Ax+By+C=0$ 旋转体体积与侧面积的通用微分积分模型（通用公式法与投影变形补偿法则）
\end{itemize}
\end{targetbox}

\begin{targetbox}[本章复习目标与核心知识图谱]
\begin{itemize}
  \item 目的[2]:深刻掌握古尔丁第一定理（Pappus-Guldin 弧线转动侧面积定理）与古尔丁第二定理（平面转动体积定理），实现形心与旋转几何量秒算互逆
\end{itemize}
\end{targetbox}

\begin{targetbox}[本章复习目标与核心知识图谱]
\begin{itemize}
  \item 目的[3]:攻克极坐标曲线绕极轴及过原点任意射线的旋转体体积与侧面积公式推导与极简积分
\end{itemize}
\end{targetbox}

\begin{targetbox}[本章复习目标与核心知识图谱]
\begin{itemize}
  \item 目的[4]:掌握已知复杂倾斜平行截面面积（椭圆截面、切角棱柱）的柱体与锥体体积高阶切片积分
\end{itemize}
\end{targetbox}

\begin{targetbox}[本章复习目标与核心知识图谱]
\begin{itemize}
  \item 目的[5]:深刻掌握无穷区间反常几何量（反常面积、反常旋转体体积、反常弧长）与级数敛散性交织的综合大题
\end{itemize}
\end{targetbox}

\begin{targetbox}[本章复习目标与核心知识图谱]
\begin{itemize}
  \item 目的[6]:掌握形心坐标、质心转动惯量与一阶/二阶矩在变力做功和液体压力中的物理-几何交叉渗透
\end{itemize}
\end{targetbox}


\section{1.斜轴旋转体体积的通用积分公式与几何补偿本质}

\begin{thmbox}[模型[1]: 绕任意一般式直线 $Ax+By+C=0$ 旋转体积通用公式(10-1)的严格推导]

在考研高等数学的高分冲刺与深水区中，旋转轴不再局限于标准的坐标轴或平行线，而是任意一般式倾斜直线 $L: Ax + By + C = 0$ ($A^2+B^2 > 0$)。

\textbf{[1]} 正交旋转变换与几何微元导出：
设平面光滑曲线 $C: y = f(x) \quad (a \leqslant x \leqslant b)$ 与直线 $L$ 无交点。

\textbf{(1)}  \textit{}回转半径 $r(x)$\textit{}：

曲线上动点 $(x, y(x))$ 到直线 $Ax + By + C = 0$ 的点到直线垂直欧氏距离为：


\begin{equation*}
r(x) = \frac{|Ax + By(x) + C|}{\sqrt{A^2 + B^2}}
\end{equation*}


\textbf{(2)}  \textit{}轴向投影微分微元 $\mathrm{d}u$\textit{}：

取直线 $L$ 的单位方向向量 $\vec{\tau} = \left(\frac{-B}{\sqrt{A^2+B^2}}, \frac{A}{\sqrt{A^2+B^2}}\right)$。曲线位移微元向量 $\mathrm{d}\vec{r} = (\mathrm{d}x, \mathrm{d}y) = (1, y'(x))\,\mathrm{d}x$ 在直线 $L$ 上的切向投影长度为：


\begin{equation*}
\mathrm{d}u = |\mathrm{d}\vec{r} \cdot \vec{\tau}| = \frac{|-B\,\mathrm{d}x + A\,\mathrm{d}y|}{\sqrt{A^2+B^2}} = \frac{|B - A y'(x)|}{\sqrt{A^2+B^2}}\,\mathrm{d}x
\end{equation*}


\textbf{[2]} 通用斜轴旋转圆盘体积公式 (10-1)：
以直线 $L$ 为轴，垂直于 $L$ 切出的圆盘微元半径为 $r(x)$，厚度为 $\mathrm{d}u$。体积微元为：


\begin{equation*}
\mathrm{d}V = \pi r^2(x)\,\mathrm{d}u = \pi \left(\frac{Ax+By+C}{\sqrt{A^2+B^2}}\right)^2 \frac{|B - A y'(x)|}{\sqrt{A^2+B^2}}\,\mathrm{d}x
\end{equation*}


整理即得考研斜轴旋转体通用终极公式：


\begin{equation*}
V = \frac{\pi}{(A^2+B^2)^{\frac{3}{2}}} \int_a^b (Ax + By(x) + C)^2 |B - A y'(x)|\,\mathrm{d}x
\end{equation*}


\textbf{[3]} 退化特例的精确自洽性验证：
\textbf{(1)}  \textit{}绕 $x$ 轴旋转\textit{}：直线方程为 $y = 0$，即 $A = 0, B = 1, C = 0$。代入公式：


\begin{equation*}
V = \frac{\pi}{(0+1)^{3/2}} \int_a^b (y(x))^2 |1 - 0|\,\mathrm{d}x = \pi \int_a^b y^2(x)\,\mathrm{d}x
\end{equation*}


完全退化为标准圆盘法公式！

\textbf{(2)}  \textit{}绕 $y$ 轴旋转\textit{}：直线方程为 $x = 0$，即 $A = 1, B = 0, C = 0$。代入公式：


\begin{equation*}
V = \frac{\pi}{(1+0)^{3/2}} \int_a^b (x)^2 |0 - y'(x)|\,\mathrm{d}x = \pi \int_a^b x^2 |y'(x)|\,\mathrm{d}x = \pi \int_{y(a)}^{y(b)} x^2\,\mathrm{d}y
\end{equation*}


若分部积分 $\int x^2 y'\,\mathrm{d}x = [x^2 y] - 2\int x y\,\mathrm{d}x$，即与圆柱壳法完全等价！

\end{thmbox}
\begin{thmbox}[模型[2]: 倾斜圆盘微元的几何补偿与斜交投影误差对消（板书 $\alpha = \beta$ 深度剖析）]

在考研例题 10.14 中，板书特别强调了斜切微元的几何补偿机理：

\textbf{[1]} 为什么非正交切片不能直接当做圆台？
若在直角坐标系中用垂直于 $x$ 轴的垂线切出宽为 $\Delta x$ 的窄条，该窄条绕斜直线 $L$ 旋转时，其横截面与旋转轴不垂直，而是倾斜了一个夹角 $\theta$。

倾斜旋转体在两端的侧表面形成了倾斜的椭圆截面圆台，直接计算截面面积会导致“几何形状失真”。

\textbf{[2]} 高阶误差的对称抵消（$\alpha = \beta$ 补偿机制）：
当曲线光滑连续时，薄窄条的左侧倾角与右侧倾角在 $\Delta x \to 0$ 时满足：


\begin{equation*}
\alpha = \beta + O(\Delta x)
\end{equation*}


在窄条绕斜轴旋转一周的立体中，凸出圆盘的部分与凹进圆盘的部分在体积上产生了一阶等价的体积补偿：


\begin{equation*}
\Delta V_{\text{left}} - \Delta V_{\text{right}} = o(\Delta u)
\end{equation*}


因此，虽然单个倾斜切片并非严格的平底圆柱，但其体积关于轴向厚度 $\mathrm{d}u$ 的主部恰好严格等于正交圆盘 $\pi r^2\,\mathrm{d}u$。这正是微元法允许自由选择投影方向的深层几何根基。

\end{thmbox}

\section{2.古尔丁定理(Pappus-Guldin)与形心坐标的逆向秒杀}

古尔丁定理（又称巴普斯定理，Pappus-Guldin Theorem）是连接几何图形面积、弧长与旋转体体积、旋转侧面积的一座宏伟大桥。

\begin{thmbox}[模型[3]: 古尔丁第一定理（旋转曲面侧面积定理）]

\textbf{[1]} 定理表述：
一条平面光滑曲线 $L$ 绕同一平面内一条与该曲线不相交的直线轴旋转一周所得旋转曲面的侧面积 $S$，等于该曲线的弧长 $s$ 乘以曲线形心（质心）绕旋转轴旋转一周所经过的圆周路程 $2\pi d_C$：


\begin{equation*}
S = 2\pi d_C \cdot s
\end{equation*}


其中 $d_C$ 为曲线形心到旋转轴的垂直欧氏距离。

\textbf{[2]} 定理证明：
设曲线形心到旋转轴的距离为 $d(x, y)$。根据曲线质心定义：


\begin{equation*}
d_C = \frac{\int_L d(x, y)\,\mathrm{d}s}{\int_L \mathrm{d}s} = \frac{1}{s}\int_L d(x, y)\,\mathrm{d}s
\end{equation*}


两端同乘 $2\pi s$：


\begin{equation*}
2\pi d_C \cdot s = 2\pi \int_L d(x, y)\,\mathrm{d}s
\end{equation*}


而根据旋转侧面积微元定义，$\mathrm{d}S = 2\pi d(x, y)\,\mathrm{d}s$，其积分恰好为旋转曲面的侧面积 $S$！定理得证。

\end{thmbox}
\begin{thmbox}[模型[4]: 古尔丁第二定理（旋转体体积定理）]

\textbf{[1]} 定理表述：
一个平面闭合区域 $D$ 绕同一平面内一条与该区域不相交（至多在边界上相切）的直线轴旋转一周所得旋转体的体积 $V$，等于该区域的面积 $A$ 乘以区域形心绕旋转轴旋转一周所经过的圆周路程 $2\pi d_C$：


\begin{equation*}
V = 2\pi d_C \cdot A
\end{equation*}


其中 $d_C$ 为平面区域形心到旋转轴的垂直欧氏距离。

\textbf{[2]} 考研秒杀应用：
\textbf{(1)}  \textbf{圆环体（Torus）体积秒杀}：

圆盘 $x^2 + (y-b)^2 \leqslant k^2$ ($0 < k < b$) 绕 $x$ 轴旋转：

- 圆盘面积 $A = \pi k^2$；
- 形心即圆心 $(0, b)$，到 $x$ 轴距离 $d_C = b$；
- 秒得体积：$V = 2\pi b \cdot (\pi k^2) = 2\pi^2 k^2 b$。（瞬间秒杀 Homework 10.2！）

\textbf{(2)}  \textbf{半圆盘形心坐标反求}：

上半圆盘 $x^2+y^2 \leqslant R^2$ ($y \geqslant 0$) 绕 $x$ 轴旋转一周得球体，其体积为 $\frac{4}{3}\pi R^3$。

半圆盘面积为 $A = \frac{1}{2}\pi R^2$。由古尔丁定理：


\begin{equation*}
V = 2\pi \overline{y} \cdot A \implies \frac{4}{3}\pi R^3 = 2\pi \overline{y} \left(\frac{1}{2}\pi R^2\right) = \pi^2 R^2 \overline{y}
\end{equation*}


秒解半圆盘形心纵坐标：


\begin{equation*}
\overline{y} = \frac{4R}{3\pi}
\end{equation*}


\end{thmbox}

\section{3.极坐标系下的旋转体体积与侧面积}

\begin{thmbox}[模型[5]: 极坐标图形绕极轴与垂直射线旋转的体积微元公式]

\textbf{[1]} 绕极轴（$x$ 轴）旋转的球台扇形微元：
设平面图形由极坐标射线 $\theta = \alpha, \theta = \beta$ 及曲线 $r = r(\theta)$ 围成。

取夹角为 $\mathrm{d}\theta$ 的扇形薄片，其面积为 $\frac{1}{2}r^2\,\mathrm{d}\theta$。扇形的重心到原点的距离为 $\frac{2}{3}r$。

重心的纵坐标（到极轴的距离）为 $y_C = \frac{2}{3}r\sin\theta$。

根据古尔丁第二定理，该极坐标小扇形绕极轴旋转一周生成的立体体积微元为：


\begin{equation*}
\mathrm{d}V_x = 2\pi y_C \cdot \mathrm{d}A = 2\pi \left(\frac{2}{3}r\sin\theta\right) \left(\frac{1}{2}r^2\,\mathrm{d}\theta\right) = \frac{2}{3}\pi r^3(\theta)\sin\theta\,\mathrm{d}\theta
\end{equation*}


故极坐标绕极轴旋转的总体积公式为：


\begin{equation*}
V_x = \frac{2}{3}\pi \int_\alpha^\beta r^3(\theta)\sin\theta\,\mathrm{d}\theta
\end{equation*}


\textbf{[2]} 绕垂直射线（$\theta = \frac{\pi}{2}$，即 $y$ 轴）旋转：
重心到 $y$ 轴的距离为 $x_C = \frac{2}{3}r\cos\theta$。


\begin{equation*}
\mathrm{d}V_y = \frac{2}{3}\pi r^3(\theta)\cos\theta\,\mathrm{d}\theta \implies V_y = \frac{2}{3}\pi \int_\alpha^\beta r^3(\theta)\cos\theta\,\mathrm{d}\theta
\end{equation*}


\textbf{[3]} 经典心形线体积秒杀：
心形线 $r = a(1+\cos\theta)$ ($0 \leqslant \theta \leqslant \pi$) 绕极轴旋转所得立体的体积：


\begin{equation*}
V_x = \frac{2}{3}\pi a^3 \int_0^\pi (1+\cos\theta)^3 \sin\theta\,\mathrm{d}\theta = \frac{2}{3}\pi a^3 \int_0^\pi (1+\cos\theta)^3 (-\mathrm{d}(1+\cos\theta))
\end{equation*}



\begin{equation*}
= \frac{2}{3}\pi a^3 \left[-\frac{(1+\cos\theta)^4}{4}\right]_0^\pi = \frac{2}{3}\pi a^3 \left(0 - \left(-\frac{16}{4}\right)\right) = \frac{8}{3}\pi a^3
\end{equation*}


\end{thmbox}

\begin{equation*}
S_x = 2\pi \int_\alpha^\beta r(\theta)\sin\theta \sqrt{r^2(\theta) + [r'(\theta)]^2}\,\mathrm{d}\theta
\end{equation*}
\begin{thmbox}[模型[6]: 极坐标曲线绕极轴旋转的侧面积微元公式]

曲线上点到极轴距离为 $y = r(\theta)\sin\theta$，极坐标弧长微元为 $\mathrm{d}s = \sqrt{r^2(\theta) + [r'(\theta)]^2}\,\mathrm{d}\theta$。


\end{thmbox}

\section{4.已知平行截面立体体积的高阶切片模型}

\begin{thmbox}[模型[7]: 双圆柱直交相交体（牟合方盖 Steinmetz Solid）的体积与侧面积]

\textbf{[1]} 问题背景：
两个底半径均为 $R$ 的直圆柱面 $x^2 + y^2 \leqslant R^2$ 与 $x^2 + z^2 \leqslant R^2$ 正交相交，求相交公共立体的体积。

\textbf{[2]} 平行切片法推导：
取垂直于 $x$ 轴的平面截该立体。对于给定的 $x \in [-R, R]$：

- 由 $x^2 + y^2 \leqslant R^2$ 可得：$-\sqrt{R^2-x^2} \leqslant y \leqslant \sqrt{R^2-x^2}$，截面在 $y$ 方向跨度为 $2\sqrt{R^2-x^2}$；
- 由 $x^2 + z^2 \leqslant R^2$ 可得：$-\sqrt{R^2-x^2} \leqslant z \leqslant \sqrt{R^2-x^2}$，截面在 $z$ 方向跨度为 $2\sqrt{R^2-x^2}$。

因此，在垂直于 $x$ 轴的截面处，形状为一个边长为 $2\sqrt{R^2-x^2}$ 的\textbf{正方形}！

截面面积为：


\begin{equation*}
A(x) = \left(2\sqrt{R^2-x^2}\right)^2 = 4(R^2 - x^2)
\end{equation*}


\textbf{[3]} 定积分计算总体积：

\begin{equation*}
V = \int_{-R}^R A(x)\,\mathrm{d}x = 8 \int_0^R (R^2 - x^2)\,\mathrm{d}x = 8 \left[R^2 x - \frac{x^3}{3}\right]_0^R = 8 \left(R^3 - \frac{R^3}{3}\right) = \frac{16}{3}R^3
\end{equation*}


\end{thmbox}
\begin{thmbox}[模型[8]: 倾斜截角圆柱体与水槽切面问题（Cavalieri 楔体）]

底圆半径为 $R$ 的圆柱体被一个过底圆直径且与底面夹角为 $\alpha$ ($0 < \alpha < \frac{\pi}{2}$) 的平面所截，求底面与截面之间所夹的楔形体体积。

\textbf{[1]} 建立坐标系与截面微元：
以底圆直径为 $x$ 轴，垂直于直径的半径方向为 $y$ 轴。底圆方程为 $x^2 + y^2 \leqslant R^2$ ($y \geqslant 0$)。

取垂直于 $x$ 轴的截面，在任意 $x \in [-R, R]$ 处：

- 底边长为 $y = \sqrt{R^2 - x^2}$；
- 截面为直角三角形，倾角为 $\alpha$，直角高为 $h = y\tan\alpha = \sqrt{R^2-x^2}\tan\alpha$。

截面三角形面积为：


\begin{equation*}
A(x) = \frac{1}{2} y \cdot h = \frac{1}{2}(R^2 - x^2)\tan\alpha
\end{equation*}


\textbf{[2]} 定积分计算总体积：

\begin{equation*}
V = \int_{-R}^R A(x)\,\mathrm{d}x = \tan\alpha \int_0^R (R^2 - x^2)\,\mathrm{d}x = \tan\alpha \left(R^3 - \frac{R^3}{3}\right) = \frac{2}{3}R^3\tan\alpha
\end{equation*}


\end{thmbox}

\section{5.反常积分与无穷几何量级数敛散性联动模型}

\begin{thmbox}[模型[9]: 托里拆利小号（Gabriel's Horn）悖论与反常几何量的 $p$ 级数本质]

将曲线 $y = \frac{1}{x} \quad (x \geqslant 1)$ 绕 $x$ 轴旋转一周所得的旋转曲面及立体称为“托里拆利小号”。

\textbf{[1]} 旋转体体积计算（收敛）：

\begin{equation*}
V = \pi \int_1^{+\infty} y^2(x)\,\mathrm{d}x = \pi \int_1^{+\infty} \frac{1}{x^2}\,\mathrm{d}x = \pi \left[-\frac{1}{x}\right]_1^{+\infty} = \pi
\end{equation*}


体积为一个完全收敛的有限常数 $\pi$！

\textbf{[2]} 旋转曲面表面积计算（发散）：

\begin{equation*}
S = 2\pi \int_1^{+\infty} y\sqrt{1+(y')^2}\,\mathrm{d}x = 2\pi \int_1^{+\infty} \frac{1}{x}\sqrt{1 + \frac{1}{x^4}}\,\mathrm{d}x
\end{equation*}


由于对任意 $x \geqslant 1$，$\sqrt{1 + \frac{1}{x^4}} > 1$：


\begin{equation*}
S > 2\pi \int_1^{+\infty} \frac{1}{x}\,\mathrm{d}x = 2\pi [\ln x]_1^{+\infty} = +\infty
\end{equation*}


表面积严格发散至无穷大！

\textbf{[3]} 数理本质深刻剖析：
“可以用有限的油漆填满小号的内部，却无法用这桶油漆刷满小号的内表面”——这表面上的物理悖论，在微积分中的本质是无穷反常积分的 $p$ 级数敛散性界限：

- 体积微元是被积函数的平方，对应 $p = 2 > 1$，反常积分收敛；
- 侧面积微元是被积函数的一次方，对应 $p = 1 \leqslant 1$，反常积分发散。

考研数学以此模型考察学生对反常积分审敛法及几何微元本质的融会贯通。

\end{thmbox}
\begin{thmbox}[模型[10]: 交错环瓣无限面积级数与双曲函数归纳（通式导出）]

设曲线 $y = e^{-kx}\sin(\omega x)$ ($k > 0, \omega > 0, x \geqslant 0$)，求该曲线与 $x$ 轴围成的全部叶瓣面积之和。

\textbf{[1]} 单叶面积积分：
交点为 $x_n = \frac{n\pi}{\omega}$ ($n = 0, 1, 2, \dots$)。

第 $n$ 个叶瓣的面积为：


\begin{equation*}
S_n = \int_{\frac{n\pi}{\omega}}^{\frac{(n+1)\pi}{\omega}} |e^{-kx}\sin(\omega x)|\,\mathrm{d}x = e^{-\frac{nk\pi}{\omega}} \int_0^{\frac{\pi}{\omega}} e^{-kt}\sin(\omega t)\,\mathrm{d}t
\end{equation*}


分部积分计算基本项：


\begin{equation*}
\int_0^{\frac{\pi}{\omega}} e^{-kt}\sin(\omega t)\,\mathrm{d}t = \frac{\omega}{k^2+\omega^2}\left(1 + e^{-\frac{k\pi}{\omega}}\right)
\end{equation*}


\textbf{[2]} 等比级数求和与双曲函数统一：
公比为 $q = e^{-\frac{k\pi}{\omega}} < 1$。全部面积之和为：


\begin{equation*}
S = \sum_{n=0}^\infty S_n = \frac{\omega}{k^2+\omega^2}\left(1 + e^{-\frac{k\pi}{\omega}}\right) \frac{1}{1 - e^{-\frac{k\pi}{\omega}}} = \frac{\omega}{k^2+\omega^2} \coth\left(\frac{k\pi}{2\omega}\right)
\end{equation*}


当 $k=1, \omega=1$ 时，立即重现板书例 10.4 的结果：$S = \frac{1}{2}\coth\frac{\pi}{2} = \frac{e^\pi+1}{2(e^\pi-1)}$。通式与特例完美统一！
\end{thmbox}

\end{document}
